% Round 3 geometry/discovery fragment. No preamble; all results proved below.
% Environments theorem, proposition, lemma, counterexample, proof assumed.
% Labels use the prefix rthreegeo to avoid collisions.

\section{Geometric Bayes risk and explicit regularity bridges}

\begin{proposition}[Coherent geometric value of information]
\label{rthreegeo:coherence}
Let $G$ be a random geometric target, $\mathcal A$ a fixed action space,
and $L(a,G)\ge0$ an integrable loss. Given data $D$, define
\[
 R(D)=\inf_{a\in\mathcal A}\E[L(a,G)\mid D],\qquad
 \Delta_j(D)=R(D)-\E[R(D,O_j)\mid D].
\]
Under coherent conditioning, $\Delta_j(D)\ge0$ whenever the conditional
Bayes risks and decisions are measurable. This is an expected, one-step
statement, not an outcome-wise or multistep-optimality statement.
\end{proposition}
\begin{proof}
Retain a current Bayes action $a_D$ after every possible observation.
Optimality of the updated action gives
$R(D,O_j)\le\E[L(a_D,G)\mid D,O_j]$.
Take the conditional expectation and use the tower identity. If the
infimum is not attained, apply this argument to measurable
$\epsilon$-optimal actions and let $\epsilon\downarrow0$.
\end{proof}

\begin{proposition}[Graph volume, medians, means and binary observations]
\label{rthreegeo:graphloss}
Let $U\subset\mathbb R^m$, let $g,h:U\to[a,b]$ be measurable, and put
$A_g=\{(u,z):u\in U,\ a\le z\le g(u)\}$.
For a nonnegative density $q$,
\[
 \mu_q(A_g\triangle A_h)
 =\int_U\left|\int_{g(u)}^{h(u)}q(u,z)\,dz\right|du.
\]
If $q(u,z)=w(u)$, this is $\int_Uw(u)|g(u)-h(u)|du$.
If $0<c\le q\le C$ between the graphs, it lies between
$c\|g-h\|_1$ and $C\|g-h\|_1$.

For a random graph $g$, the Bayes action under integrated absolute
displacement is a measurable pointwise posterior median. Under
integrated squared displacement, it is the posterior mean, and the
exact value of observing $O_j$ is
\[
 \Delta_{2,j}
 =\int_Uw(u)\operatorname{Var}\bigl(\E[g(u)\mid D,O_j]\mid D\bigr)\,du.
\]
For a binary observation $O_j\in\{-1,1\}$ with positive variance,
\[
 \Delta_{2,j}
 =\int_Uw(u)
 \frac{\operatorname{Cov}(g(u),O_j\mid D)^2}
      {\operatorname{Var}(O_j\mid D)}\,du.
\]
If every posterior graph is $L$-Lipschitz, the pointwise posterior mean
and a consistently chosen quantile median are also $L$-Lipschitz.
\end{proposition}
\begin{proof}
At each $u$, the symmetric difference of the vertical sections is exactly
the interval between $g(u)$ and $h(u)$. Tonelli's theorem proves the
identity; the density bounds follow pointwise.
Tonelli also separates each integrated posterior loss by $u$.
For a scalar $Z$, a number $a$ minimizes $\E|a-Z|$ precisely when
$\mathbb P(Z\le a)\ge1/2$ and $\mathbb P(Z\ge a)\ge1/2$; this follows
from the one-sided derivatives of the convex objective. The squared-loss
claim follows by completing the square.
Conditional total variance proves the displayed value identity.
Every square-integrable function of a nondegenerate binary variable is
affine in that variable; its conditional-mean regression coefficient is
$\operatorname{Cov}(g(u),O_j)/\operatorname{Var}(O_j)$, giving the final
formula. A deterministic observation has value zero.
Finally, the almost-sure inequality
$|g(u)-g(v)|\le L\|u-v\|$ implies the same inequality for corresponding
quantiles of the two marginal laws. Taking expectations gives it for
the means. These actions are measurable under the usual measurable
posterior-process assumptions.
\end{proof}

\begin{proposition}[A restricted Gaussian graph problem]
\label{rthreegeo:gaussiangraph}
Suppose the current posterior of $g$ is Gaussian, with variance $s(u)^2$
and covariance $c(u,v)$, and a paid observation directly reveals $g(v)$.
Assume the displayed integrals are finite. For integrated absolute
displacement of unrestricted graph heights,
\[
 R_1=\sqrt{\frac2\pi}\int_U w(u)s(u)\,du,
\]
\[
 \Delta_1(v)=\sqrt{\frac2\pi}\int_Uw(u)
 \left[s(u)-\sqrt{s(u)^2-\frac{c(u,v)^2}{s(v)^2}}\right]du.
\]
For integrated squared displacement,
\[
 \Delta_2(v)=\int_Uw(u)\frac{c(u,v)^2}{s(v)^2}\,du.
\]
The values are zero if $s(v)=0$. A binary membership observation does not
directly reveal $g(v)$ and is a different observation experiment.
\end{proposition}
\begin{proof}
A Gaussian marginal's mean is a median and its mean absolute centered
deviation is $\sqrt{2/\pi}s(u)$. Gaussian conditioning changes the variance
to $s(u)^2-c(u,v)^2/s(v)^2$, independently of the observed value.
Substitute this into the absolute-deviation risk. For squared loss,
subtract the integrated updated variance from the integrated current
variance. A zero-variance observation is already known.
\end{proof}

\begin{proposition}[Graph Hausdorff and density-dependent $L^1$ control]
\label{rthreegeo:holderbridge}
Let $g,h$ be $L$-Lipschitz on the same compact $U$, and let
$\Gamma_g=\{(u,g(u)):u\in U\}$. Then
\[
 \frac{\|g-h\|_\infty}{\sqrt{1+L^2}}
 \le H(\Gamma_g,\Gamma_h)\le\|g-h\|_\infty.
\]
More specifically, let $U=[0,1]^m$, let $f=g-h$ be $K$-Lipschitz,
$K>0$, and let $w\ge w_{\min}>0$. Put $M=\|f\|_\infty$ and
$r=\min\{1/2,M/(2K\sqrt m)\}$. Then
\[
 \int_Uw|f|\ge w_{\min}\frac M2 r^m.
\]
In the regime $M\le K\sqrt m$ this gives
\[
 M\le\left[
 \frac{2^{m+1}K^m m^{m/2}}{w_{\min}}\int_Uw|f|
 \right]^{1/(m+1)}.
\]
\end{proposition}
\begin{proof}
Match graph points with the same base coordinate for the upper bound.
At $u$ maximizing $|g(u)-h(u)|=M$, a point $(v,h(v))$ is at least
\[
 \sqrt{t^2+(M-Lt)_+^2},\qquad t=\|u-v\|,
\]
away. The minimum over $t\ge0$ is $M/\sqrt{1+L^2}$.
For the integral bound, at a maximizer of $|f|$ choose, in each coordinate,
a direction with at least $1/2$ available space inside the unit cube.
The resulting one-sided box of side $r$ is contained in $U$. Every point
in it lies within $\sqrt m r$ of the maximizer, and therefore has
$|f|\ge M/2$. Integrate on this box and rearrange. For $K=0$, $f$ is
constant and the integral is exactly $M\int w$. A cone of height $M$
and slope $K$ has mass of order $M^{m+1}/K^m$, showing that the
small-error exponent cannot generally be improved under only a
Lipschitz bound.
\end{proof}

\begin{proposition}[A finite graph approximation]
\label{rthreegeo:quadrature}
Let a probability measure $\nu$ on $U$ be partitioned into cells $C_i$
of diameter at most $h$, with $u_i\in C_i$ and $w_i=\nu(C_i)$.
For graphs with Lipschitz constants $L_g,L_h$,
\[
 \left|\int_U|g-h|\,d\nu-
 \sum_iw_i|g(u_i)-h(u_i)|\right|\le(L_g+L_h)h.
\]
If the base points form an $h$-net, the difference between the Hausdorff
distance of the sampled graph clouds and that of the complete graphs
is at most
\[
 h\left(\sqrt{1+L_g^2}+\sqrt{1+L_h^2}\right).
\]
\end{proposition}
\begin{proof}
The function $|g-h|$ is $(L_g+L_h)$-Lipschitz. Integrate its deviation
from each cell representative. Every graph point is within
$h\sqrt{1+L_g^2}$ of its sampled graph cloud, and conversely every cloud
point lies on the graph. Apply the triangle inequality for Hausdorff
distance twice. These arguments require graph heights; a finite label
pool alone need not locate an intervening root.
\end{proof}

\begin{proposition}[A density-weighted normal-displacement expansion]
\label{rthreegeo:tube}
Let $S$ be a compact oriented $C^2$ hypersurface of dimension $m$ with
injective normal coordinates $F(s,t)=s+tn(s)$ for $|t|<r$.
Suppose its principal curvatures are bounded in magnitude by $K$.
Let $0<a<r$, $aK<1$, and assume two regions differ exactly by the
normal strip between $t=0$ and $t=\delta(s)$, where $|\delta|\le a$;
there are no additional components or orientation changes.
Let $q\ge0$ be bounded by $Q$ and $L_q$-Lipschitz on the tube. Then
\[
 \mu_q(A\triangle B)=
 \int_S\left|\int_0^{\delta(s)}q(F(s,t))J(s,t)\,dt\right|dA(s),
 \quad J(s,t)=\prod_{i=1}^m(1-t\kappa_i(s)),
\]
and
\[
 \left|\mu_q(A\triangle B)-\int_Sq(s)|\delta(s)|\,dA(s)\right|
 \le\frac C2\int_S\delta(s)^2\,dA(s),
\]
where
\[
 C=L_q(1+aK)^m+QmK(1+aK)^{m-1}.
\]
\end{proposition}
\begin{proof}
The normal-coordinate change of variables gives the first identity.
The assumptions make its Jacobian positive. Telescoping the product
gives
$|J-1|\le mK|t|(1+aK)^{m-1}$ and $|J|\le(1+aK)^m$.
Consequently
$|q(F(s,t))J(s,t)-q(s)|\le C|t|$.
Integrating from $0$ to $\delta(s)$ and then over $S$ proves the bound.
\end{proof}

\begin{counterexample}[Latent amplitude does not identify boundary uncertainty]
\label{rthreegeo:scale}
Let $\Theta\sim N(m,\sigma^2)$ and $f_c(x)=c(x-\Theta)$, $c>0$.
The boundary is $\{\Theta\}$ for all $c$, although
$\operatorname{sd}(f_c(x))=c\sigma\to0$ as $c\downarrow0$.
Squared boundary Bayes risk remains $\sigma^2$; the scale-invariant
quantity $\operatorname{sd}(f_c)/|f_c'|$ equals $\sigma$.
A noiseless sign query at $x=m$ has squared-boundary value
$2\sigma^2/\pi$. Under the different observation contract
$O=\operatorname{sign}(f_c(m)+\eta)$ with independent
$\eta\sim N(0,\tau^2)$, the value is
\[
 \frac{2c^2\sigma^4}{\pi(c^2\sigma^2+\tau^2)}.
\]
\end{counterexample}
\begin{proof}
The root and scale claims follow directly from the definition.
A noiseless query splits the Gaussian threshold at its median; the two
conditional means are $m\pm\sigma\sqrt{2/\pi}$, whose variance is
$2\sigma^2/\pi$. More generally put $X=\Theta-m$ and
$Z=-cX+\eta$. Joint Gaussian regression gives
\[
 \E[X\mid\operatorname{sign}Z]
 =\frac{\operatorname{Cov}(X,Z)}{\operatorname{sd}(Z)}
   \sqrt{\frac2\pi}\operatorname{sign}Z.
\]
Its variance is the displayed value by conditional total variance.
The fixed-noise value tends to zero while the random boundary law stays
unchanged.
\end{proof}

\begin{proposition}[Local root sensitivity requires a nonzero gradient]
\label{rthreegeo:rootsensitivity}
On a fixed normal line, let $f_t(x)=f(x)+t e(x)$ be continuously
differentiable near $(t,x)=(0,x_0)$, with $f(x_0)=0$ and
$f'(x_0)\ne0$. There is a local differentiable root $x(t)$, and
\[
 x'(0)=-\frac{e(x_0)}{f'(x_0)}.
\]
This derivative is invariant when both $f$ and $e$ are multiplied by a
positive constant.
\end{proposition}
\begin{proof}
The implicit function theorem gives the root and its differentiability.
Differentiate $f(x(t))+t e(x(t))=0$ at zero and solve for $x'(0)$.
The multiplicative constant cancels. Uniform positional bounds require
a uniform derivative lower bound and control of additional roots;
the local identity by itself supplies neither.
\end{proof}

\begin{counterexample}[Area and tolerance overlap are different losses]
\label{rthreegeo:areansd}
For $h=0$ and $g_k(u)=k^{-1}\sin(k^2u)$ on $[0,1]$, graph Hausdorff
distance and symmetric-difference area tend to zero, whereas graph
length diverges. For two parallel equal-area surfaces separated by $d$,
normalized surface Dice at tolerance $\tau$ is one for $d\le\tau$ and
zero for $d>\tau$, while Hausdorff distance is $d$.
\end{counterexample}
\begin{proof}
Uniform graph displacement is at most $1/k$, bounding both Hausdorff
distance and integrated displacement. Graph length is at least
$\int_0^1 k|\cos(k^2u)|du=(2/\pi)k+o(k)$.
For the parallel surfaces every point's nearest-surface distance is
exactly $d$, so every point either counts within tolerance or does not.
\end{proof}

\section{A Hausdorff acquisition counterexample beyond pairwise moments}

\begin{lemma}[Bayes sup-distance risk on a binary cube]
\label{rthreegeo:cuberisk}
For any distribution of $B\in\{0,1\}^m$, allowing all estimates
$h\in[0,1]^m$,
\[
 \inf_h\E\|h-B\|_\infty
 =\min\left\{\frac12,1-\max_b\mathbb P(B=b)\right\}.
\]
\end{lemma}
\begin{proof}
The cube center has risk $1/2$, and a modal vertex has risk $1-p_{\max}$.
For any $h$, round each coordinate to a nearest endpoint to obtain $v$,
and let $r=\|h-v\|_\infty\le1/2$.
Every other vertex is at distance at least $1-r$, hence the risk is at
least $p_vr+(1-p_v)(1-r)$. Its minimum over $r\in[0,1/2]$ is
$\min\{1/2,1-p_v\}$, which is at least the stated lower bound.
The two explicit actions attain that bound.
\end{proof}

\begin{counterexample}[Equal binary pair moments, different geometric query values]
\label{rthreegeo:parity}
Let three bases $u_i$ have pairwise distances at least $3$.
A random boundary is $\{(u_i,B_i):i=1,2,3\}$, where
$B\in\{0,1\}^3$, and admissible estimates have heights in $[0,1]^3$.
Querying column $i$ at height $1/2$ reveals $B_i$.
The following two laws have identical first and pairwise binary moments:
\[
\begin{array}{c|rrrrrrrr}
 B&000&001&010&011&100&101&110&111\\\hline
 70P_+(B)&34&0&0&12&0&12&12&0\\
 70P_-(B)&22&12&12&0&12&0&0&12
\end{array}
\]
Their current Hausdorff Bayes risks are both $1/2$, but the exact values
of querying $B_1$ are $11/70$ and zero, respectively.
\end{counterexample}
\begin{proof}
Because other columns are at distance at least $3$ whereas the
same-column height difference is at most $1$, the Hausdorff loss is
exactly $\|h-B\|_\infty$.
Direct summation gives, for both laws,
\[
 \mathbb P(B_i=1)=24/70,\qquad
 \mathbb P(B_i=B_j=1)=12/70\quad(i\ne j).
\]
These determine every first and pairwise binary moment.
Every atom has probability less than $1/2$, so
Lemma~\ref{rthreegeo:cuberisk} gives current risk $35/70$.
The events $B_1=0$ and $B_1=1$ have masses $46/70$ and $24/70$.
For a branch of mass $s$ and largest atom mass $p$, its unnormalized
Bayes risk is $\min\{s/2,s-p\}$ by the same lemma.
Under $P_+$, branch zero has largest atom $34/70$ and risk $12/70$;
branch one has risk $12/70$. Updated expected risk is $24/70$.
Under $P_-$, branch zero has largest atom $22/70$, so its risk is
$23/70$; branch one again has risk $12/70$. Updated expected risk
is $35/70$. Subtraction proves the two query values.
This counterexample permits intermediate height estimates and thus
does not rely on forcing the estimate to be a posterior sample.
\end{proof}

\section{Discovery, refinement, and violations of structural assumptions}

\begin{counterexample}[Optimal startup, suboptimal later geometry at equal budget]
\label{rthreegeo:discoverygeometry}
Fix any $k\ge1$, put $M=k+1$, and let
$\Theta_1,\ldots,\Theta_M$ be independent fair bits.
The geometric target has $M$ separated columns, heights $H\Theta_i$,
and base separations at least $3H$.
The finite candidate pool contains one membership query $q_i$ at height
$H/2$ in each column, revealing $\Theta_i$, and two paid certification
anchors of fixed labels $0$ and $1$.
Startup requires actually observed labels of both classes.

Policy A queries the anchors first. It has the globally minimal
$T_{\rm both}=2$ surely. With an optimal continuation and total paid
budget $Q=k+2$, its final Hausdorff Bayes risk is $H/2$.
Policy B has $\E T_{\rm both}=5/2$, $T_{\rm both}\le3$, and zero final
geometric risk by the same total budget $Q$.
\end{counterexample}
\begin{proof}
No policy can observe both labels in fewer than two queries, proving
A's startup optimality. After its two anchors, A has at most $k$
individual-bit queries for $M=k+1$ independent bits.
Conditional on any adaptive query history, an unqueried bit remains
independent and fair. Its minimum expected absolute height error is
$H/2$, a lower bound on expected Hausdorff error.
Set every unknown height to $H/2$ and every known height correctly;
the loss is exactly $H/2$, attaining the bound.

B first queries $q_1,q_2$. With probability $1/2$ their labels differ,
so startup ends at two. Otherwise it queries the opposite-label anchor,
ending at three. Thus its expected startup cost is $5/2$.
It then queries the remaining $M-2=k-1$ bits. The total cost is at most
$3+(k-1)=Q$, and all boundary heights are known. Unused budget can be
padded by an irrelevant remaining query if an exactly equal paid count
is desired. Therefore its risk is zero.

This is an operational paid-certification example. If known structural
anchors may replace actual paid observations, the startup contract
changes. The result concerns the stated actual-observation contract.
\end{proof}

\begin{proposition}[State coalescence bounds later effects]
\label{rthreegeo:coupling}
Couple two policies on the same truth and continuation seed. Suppose
their complete sufficient states coincide at budget $b$ with probability
at least $1-\delta$, and they use the same subsequent transition and
prediction rules. For any later endpoint in an interval of length $L$,
\[
 |\E Z_A-\E Z_B|\le L\delta.
\]
For a weighted learning-curve endpoint with nonnegative weights
$\lambda_t$ summing to one, if $\delta_t$ bounds the probability of
noncoalescence at budget $t$, then
\[
 |\E\textstyle\sum_t\lambda_t Z_{A,t}
       -\E\sum_t\lambda_t Z_{B,t}|
 \le L\sum_t\lambda_t\delta_t.
\]
Identical query sets imply identical states only under order-invariant
fitting, equal seeds, and absence of additional history-dependent state.
\end{proposition}
\begin{proof}
Identical transitions preserve identical states inductively on the
coalescence event. Endpoint differences vanish on that event and are
bounded by $L$ on its complement. Take expectations.
Apply this argument at each budget and use the triangle inequality for
the weighted endpoint.
\end{proof}

\begin{counterexample}[Three different meanings of monotonicity]
\label{rthreegeo:monotone}
Monotone class probabilities do not permit deterministic label
propagation: independent labels with probabilities $0.49,0.51$ are
stochastically ordered but have probability $0.49^2$ of the reversed
pair $(1,0)$.

On a chain of $N$ points, the deterministic truth
$(1,0,\ldots,0,1)$ differs at one point from the monotone threshold
$(0,\ldots,0,1)$. Propagating the first observed positive upward gives
$N-2$ wrong labels. This is one contaminated point but $N-2$ violating
pairs. Even exactly one violating pair destroys zero-error propagation:
the sequence $(0,\ldots,0,1,0,1,\ldots,1)$ with adjacent $1,0$
incorrectly propagates the first of those two labels to the second.
\end{counterexample}
\begin{proof}
Independence gives $\mathbb P(Y_1=1,Y_2=0)=0.49(1-0.51)$.
For the first deterministic sequence, every middle zero is dominated
by the first positive and is therefore incorrectly removed by the
hard monotonicity rule. Each is also a distinct violating pair with
the first point. In the second sequence the adjacent reversed pair is
the only inversion; propagation across it is nonetheless wrong.
\end{proof}

\begin{proposition}[A robust and tight rank-error discovery bound]
\label{rthreegeo:rank}
In a fixed tie-broken order of $N$ items containing $r\ge1$ designated
rare items, let $V$ count common-before-rare pairs, and let $D$ be the
first rare rank. Then
\[
 D\le1+\left\lfloor\frac Vr\right\rfloor.
\]
If $0<r<N$ and AUC is computed on this strict rare-oriented ranking,
\[
 D\le1+\lfloor(1-\mathrm{AUC})(N-r)\rfloor.
\]
For a top-$M$ subset containing $s$ rare items,
\[
 V\ge(M-s)(r-s).
\]
In particular $V<Mr$ guarantees that the first $M$ positions contain
a rare item. The first-hit bound is attained when $V=kr$ by placing
$k$ common items, all $r$ rare items, then all remaining common items.
\end{proposition}
\begin{proof}
Each of the $D-1$ common items preceding the first rare precedes all
$r$ rare items. Thus $V\ge r(D-1)$.
For a strict ranking, $1-\mathrm{AUC}=V/[r(N-r)]$.
For the tail inequality, each of the $M-s$ common items inside the
top-$M$ subset precedes every one of the $r-s$ rare items outside it.
If $s=0$, this forces $V\ge Mr$.
The displayed block ordering has exactly $kr$ inversions and first-hit
rank $k+1$, proving attainment. Half-credit AUC for unresolved ties
does not fix a deterministic first-hit rank; the declared tie order
must be used.
\end{proof}
